\documentclass[aps,prl,twocolumn,showpacs,preprintnumbers,amsmath,amssymb,superscriptaddress]{revtex4-2}

\usepackage{graphicx}
\usepackage{hyperref}
\usepackage{amsmath}
\usepackage{bm}

\begin{document}

\title{Dark energy from horizon thermodynamics: a zero-parameter prediction for DESI}

\author{Charles Grant Smith Jr.}
\affiliation{Pegasus Intellectual Capital Solutions, Chicago, IL 60601, USA}

\date{\today}

\begin{abstract}
The Gibbons-Hawking temperature of the de~Sitter horizon, $T = H/(2\pi)$, 
defines a dimensionless coupling $g_0 = 1/(2\pi)$ between boundary encoding 
and bulk expansion. Combined with the requirement that the dark energy 
equation of state satisfies $w(z{=}0) = -1$ (de~Sitter thermal equilibrium), 
this uniquely determines $w(z) = -1 + (1/2\pi)\,[1 - H_0/H(z)]$ with no 
free parameters. Tested against DESI Data Release~2 baryon acoustic 
oscillation distances jointly with compressed Planck CMB constraints, this 
prediction is preferred over $\Lambda$CDM at $\Delta\chi^2 = 6.3$ ($2.5\sigma$) 
with no additional free parameters beyond the standard cosmological set 
$\{H_0, \omega_{\rm cdm}\}$. The signal grew from $1.7\sigma$ in DR1 to 
$2.5\sigma$ in DR2. The same $1/(2\pi)$ independently predicts the galaxy-scale 
acceleration anomaly $a_u = cH_0/(2\pi)$, tested against 171 SPARC rotation 
curves. The predicted $w(z)$ flattens to $-1 + 1/(2\pi) \approx -0.84$ at 
high redshift, in contrast to the CPL extrapolation $w \to -2.17$; DESI DR3 
can distinguish the two.
\end{abstract}

\maketitle

%=====================================================================
\section{Introduction}
%=====================================================================

The Dark Energy Spectroscopic Instrument (DESI) has reported a 
$3.1\sigma$ preference for time-varying dark energy over $\Lambda$CDM 
from baryon acoustic oscillation (BAO) measurements combined with the 
cosmic microwave background (CMB)~\cite{DESI_DR2}. The standard 
phenomenological description uses two free parameters, $w_0$ and $w_a$, 
in the Chevallier-Polarski-Linder (CPL) parameterization 
$w(a) = w_0 + w_a(1-a)$~\cite{Chevallier2001,Linder2003}. Best-fit 
values are $w_0 = -0.42$, $w_a = -1.75$~\cite{DESI_DR2}, indicating 
$w > -1$ today and $w < -1$ in the past.

Here we show that a specific equation of state, derived from de~Sitter 
horizon thermodynamics with no free parameters, reproduces the DESI 
anomaly. The derivation rests on one theorem (Gibbons-Hawking~\cite{GH1977}) 
and one physical requirement (thermal equilibrium at $z=0$). The resulting 
$w(z)$ is a parameter-free prediction that can be falsified by future 
data releases.

%=====================================================================
\section{Derivation}
%=====================================================================

The de~Sitter horizon radiates at the Gibbons-Hawking 
temperature~\cite{GH1977}
\begin{equation}
T_{\rm GH} = \frac{\hbar H}{2\pi k_B}\,.
\label{eq:TGH}
\end{equation}
In a holographic framework where the de~Sitter boundary is ontologically 
primary~\cite{Bousso2002,Susskind1995}, the ratio
\begin{equation}
g_0 \equiv \frac{T_{\rm GH}}{(\hbar/k_B)\,H} = \frac{1}{2\pi}
\label{eq:g0}
\end{equation}
is the dimensionless coupling between the horizon's thermal encoding 
rate and the bulk expansion rate. This is the fraction of the expansion 
rate that constitutes irreducible thermal noise on the horizon.

We require two physical conditions on the dark energy equation of state:

\noindent\textit{Condition 1.}---At $z=0$, the horizon is in thermal 
equilibrium with itself: the encoding rate matches the expansion rate 
exactly, and the dark energy equation of state takes its equilibrium 
value,
\begin{equation}
w(z{=}0) = -1\,.
\label{eq:cond1}
\end{equation}

\noindent\textit{Condition 2.}---The deviation from $w=-1$ is 
proportional to the dimensionless mismatch between the current expansion 
rate $H(z)$ and its equilibrium value $H_0$,
\begin{equation}
\delta(z) = \frac{H(z) - H_0}{H(z)} = 1 - \frac{H_0}{H(z)}\,.
\label{eq:delta}
\end{equation}
This is the unique dimensionless ratio of $H(z)$ and $H_0$ that 
(a) vanishes at $z=0$, (b) approaches unity as $z \to \infty$, 
(c) is bounded between 0 and 1, and (d) introduces no new scales.
The linear (first-order) dependence on $\delta(z)$ is mandated by 
perturbation theory: since $g_0 = 1/(2\pi) \approx 0.16 \ll 1$, 
the deviation from $w=-1$ is a small perturbation, and higher-order 
terms $\delta^n$ contribute at $\mathcal{O}(g_0^n)$, which is 
negligible at current precision.

Combining Eqs.~\eqref{eq:g0}--\eqref{eq:delta}:
\begin{equation}
\boxed{w(z) = -1 + \frac{1}{2\pi}\left[1 - \frac{H_0}{H(z)}\right]}
\label{eq:wz}
\end{equation}
with $H(z) = H_0\sqrt{\Omega_m(1+z)^3 + \Omega_\Lambda}$ as in 
$\Lambda$CDM. This prediction has no free parameters beyond 
the standard cosmological set $\{H_0, \omega_{\rm cdm}\}$.

We note that the $H(z)$ appearing in Eq.~\eqref{eq:delta} is 
strictly the $\Lambda$CDM Hubble parameter. A self-consistent 
treatment would solve for $H(z)$ iteratively; however, since 
$g_0 \approx 0.16$ and the resulting dark energy density deviates 
from $\rho_\Lambda$ by only a few percent, the error from this 
approximation is well below the statistical uncertainty of 
current data.

The dark energy density evolves as
\begin{equation}
\frac{\rho_{\rm DE}(z)}{\rho_\Lambda} = 
\exp\!\left[-3\int_0^z \frac{1+w(z')}{1+z'}\,dz'\right],
\label{eq:rhoDE}
\end{equation}
giving $\rho_{\rm DE}/\rho_\Lambda \approx 1.04$ at $z=0.7$ and 
$\approx 1.07$ at $z=1$---a 4--7\% excess dark energy density in 
the past that shortens distances relative to $\Lambda$CDM, consistent 
with the DESI measurements.

\textit{Connection to galaxy scales.}---The same $1/(2\pi)$ appears 
independently at galactic scales: the acceleration threshold 
$a_u = cH_0/(2\pi) = 1.22\times10^{-10}$~m~s$^{-2}$ marks the 
onset of the rotation curve anomaly. Tested against 171 SPARC galaxies 
with zero free parameters, $a_u$ outperforms the empirically calibrated 
MOND threshold $a_0$ in 59.6\% of cases~\cite{Paper1}. Both predictions 
originate from the Gibbons-Hawking temperature applied to different 
regimes of the same holographic framework.

%=====================================================================
\section{Data and method}
%=====================================================================

We use DESI DR2 BAO distance measurements~\cite{DESI_DR2}: 
$D_V/r_d$ at $z=0.295$ (BGS), $D_M/r_d$ and $D_H/r_d$ at 
$z = 0.510$, $0.706$, $0.934$, $1.321$, $1.484$, and $2.33$ 
(LRG, ELG, QSO, Ly$\alpha$), totaling 13 data points with the 
full $13 \times 13$ covariance matrix from 
Ref.~\cite{CobayaSampler}. The sound horizon is computed from 
the fitting formula of Ref.~\cite{Aubourg2015}.

CMB information enters as a compressed Gaussian likelihood on 
$\omega_b = 0.02237 \pm 0.00015$, 
$\omega_{\rm cdm} = 0.1200 \pm 0.0012$, and 
$H_0 r_d = 9906 \pm 23$~km~s$^{-1}$, following the approach 
described in Ref.~\cite{DESI_DR2}. No supernova data are included.

For each model, we minimize the total 
$\chi^2 = \chi^2_{\rm CMB} + \chi^2_{\rm BAO}$, where 
$\chi^2_{\rm BAO} = (\mathbf{d}-\mathbf{t})^T C^{-1}(\mathbf{d}-\mathbf{t})$ 
uses the full covariance, over the parameters $\{H_0, \omega_{\rm cdm}\}$. 
The predicted equation of state, Eq.~\eqref{eq:wz}, enters through 
Eq.~\eqref{eq:rhoDE} in the Friedmann equation, modifying $H(z)$ and 
hence all distance measures. The $\Lambda$CDM comparison uses the same 
parameter set and the same data, with $w=-1$ fixed.

%=====================================================================
\section{Results}
%=====================================================================

The best-fit $\Lambda$CDM model (2 parameters: $H_0$, $\omega_{\rm cdm}$) 
gives $\chi^2_{\Lambda\rm CDM} = 27.9$, split as 
$\chi^2_{\rm CMB} = 4.5$ and $\chi^2_{\rm BAO} = 23.4$. The nonzero 
CMB contribution reflects the known tension between Planck-preferred 
parameters and DESI-preferred distances.

The kernel prediction, Eq.~\eqref{eq:wz}, with the same 2 parameters 
gives
\begin{equation}
\chi^2_{\rm kernel} = 21.6\,,\qquad 
(\chi^2_{\rm CMB} = 1.9,\;\chi^2_{\rm BAO} = 19.7)\,.
\label{eq:chi2}
\end{equation}
The improvement is
\begin{equation}
\Delta\chi^2 = \chi^2_{\Lambda\rm CDM} - \chi^2_{\rm kernel} = 6.3
\label{eq:dchi2}
\end{equation}
for \emph{no} additional parameters beyond those already present 
in $\Lambda$CDM. The kernel simultaneously 
improves the BAO fit and reduces the CMB tension (from 
$\chi^2_{\rm CMB} = 4.5$ to $1.9$), because the modified expansion 
history allows $H_0$ and $\omega_{\rm cdm}$ to sit closer to their 
Planck-preferred values.

The prediction improves 8 of 13 individual data points, with the 
largest improvements at the redshifts where DESI's constraints are 
tightest ($z = 0.7$--$0.9$, LRG bins). The improvement is robust to 
data subset selection: removing any single data point yields 
$\Delta\chi^2 \geq 4.2$ ($\geq 2.1\sigma$), confirming that the 
signal is not driven by any individual measurement.

\textit{DR1 to DR2 evolution.}---Repeating the analysis with DESI 
DR1 data~\cite{DESI_DR1} gives $\Delta\chi^2 = 2.9$ ($1.7\sigma$). 
The signal grew to $2.5\sigma$ in DR2, consistent with the 
$\sim\!\sqrt{2}$ improvement expected from the roughly doubled 
effective volume.

%=====================================================================
\section{Falsifiable prediction}
%=====================================================================

The kernel prediction and the CPL best-fit agree at $z \approx 0.3$--$0.7$, 
where DESI's current constraints are tightest, but diverge sharply at 
higher redshift:
\begin{equation}
w_{\rm kernel}(z{\to}\infty) = -1 + \frac{1}{2\pi} \approx -0.84\,,
\end{equation}
while $w_{\rm CPL}(z{\to}\infty) = w_0 + w_a = -2.17$. At $z=2$, 
$w_{\rm kernel} = -0.89$ versus $w_{\rm CPL} = -1.59$. This 
divergence is well within the reach of DESI DR3 and the completed 
DESI survey, which will extend precise BAO measurements to $z > 2$ 
with substantially reduced errors on $D_H/r_d$.

%=====================================================================
\section{Discussion}
%=====================================================================

The result reported here---$\Delta\chi^2 = 6.3$ with zero additional 
parameters---should be compared to the DESI collaboration's finding 
of $\Delta\chi^2_{\rm MAP} = -10.7$ to $-21.0$ for $w_0 w_a$CDM 
over $\Lambda$CDM using CMB + BAO + supernovae, at 2.8--4.2$\sigma$ 
but with \emph{two} extra parameters~\cite{DESI_DR2}. The kernel model 
achieves $\Delta\chi^2 = 6.3$ at $2.5\sigma$ with no extra parameters.

The physical interpretation is that DESI's ``evolving dark energy'' 
signal is not a change in the dark energy equation of state but the 
imprint of horizon thermodynamics on the expansion history. The 
de~Sitter boundary encodes bulk degrees of freedom at a rate set by 
$T_{\rm GH}$. When $H(z) > H_0$ (as at all $z > 0$), the mismatch 
between the current and equilibrium encoding rates produces a slight 
excess in the effective dark energy density, shortening distances 
relative to $\Lambda$CDM. This is the same mechanism that, at 
galactic scales, produces the acceleration anomaly below 
$a_u = cH_0/(2\pi)$~\cite{Paper1}.

Two methodological caveats deserve emphasis. First, the compressed 
CMB likelihood is an approximation to the full Planck analysis; a 
complete treatment using CAMB within the Cobaya framework would 
provide a definitive significance. Second, the derivation of 
Eq.~\eqref{eq:wz} proceeds from dimensional analysis and equilibrium 
requirements; a microphysical derivation from the boundary's quantum 
information dynamics remains future work.

\textit{Summary.}---One number, $1/(2\pi)$, derived from the 
Gibbons-Hawking temperature, predicts both the galaxy-scale 
acceleration anomaly and the DESI cosmological-scale dark energy 
anomaly, with no additional free parameters at either scale. The prediction 
$w(z) = -1 + (1/2\pi)[1-H_0/H(z)]$ is falsifiable: it flattens to 
$w = -0.84$ at high redshift, while the CPL parameterization predicts 
$w = -2.17$. DESI DR3 will distinguish the two.

%=====================================================================
\begin{acknowledgments}
The author thanks Anthropic's Claude for computational assistance 
in data analysis and numerical verification. This work uses public 
DESI DR2 BAO data from Ref.~\cite{CobayaSampler} and Planck 2018 
compressed CMB constraints. DESI is managed by Lawrence Berkeley 
National Laboratory under U.S. Department of Energy Contract 
No.\ DE-AC02-05CH11231.
\end{acknowledgments}

\begin{thebibliography}{99}

\bibitem{DESI_DR2}
DESI Collaboration,
``DESI DR2 Results II: Measurements of Baryon Acoustic Oscillations and 
Cosmological Constraints,''
arXiv:2503.14738 [astro-ph.CO] (2025).

\bibitem{Chevallier2001}
M.~Chevallier and D.~Polarski,
Int.\ J.\ Mod.\ Phys.\ D \textbf{10}, 213 (2001).

\bibitem{Linder2003}
E.~V.~Linder,
Phys.\ Rev.\ Lett.\ \textbf{90}, 091301 (2003).

\bibitem{GH1977}
G.~W.~Gibbons and S.~W.~Hawking,
Phys.\ Rev.\ D \textbf{15}, 2738 (1977).

\bibitem{Bousso2002}
R.~Bousso,
Rev.\ Mod.\ Phys.\ \textbf{74}, 825 (2002).

\bibitem{Susskind1995}
L.~Susskind,
J.\ Math.\ Phys.\ \textbf{36}, 6377 (1995).

\bibitem{Paper1}
C.~G.~Smith~Jr.,
``Holographic decoherence threshold in SPARC galaxy rotation curves,''
submitted to Phys.\ Rev.\ Lett.\ (2026).

\bibitem{CobayaSampler}
DESI BAO likelihood data,
\url{https://github.com/CobayaSampler/bao_data}.

\bibitem{Aubourg2015}
\'E.~Aubourg \textit{et al.},
Phys.\ Rev.\ D \textbf{92}, 123516 (2015).

\bibitem{DESI_DR1}
DESI Collaboration,
``DESI 2024 VI: Cosmological Constraints from Baryon Acoustic 
Oscillations,''
arXiv:2404.03002 [astro-ph.CO] (2024).

\bibitem{DESI_KP5}
DESI Collaboration,
``DESI 2024 V: Full-Shape Galaxy Clustering from Galaxies and 
Quasars,''
arXiv:2411.12021 [astro-ph.CO] (2024).

\bibitem{DESI_KP7B}
DESI Collaboration,
``DESI 2024 VII: Cosmological Constraints from the Full-Shape 
Modeling of Clustering Measurements,''
arXiv:2411.12022 [astro-ph.CO] (2024).

\end{thebibliography}

\end{document}
